\documentclass[11pt,a4paper]{article}
\usepackage[utf8]{inputenc}
\usepackage[T1]{fontenc}
\usepackage[turkish]{babel}
\usepackage{geometry}
\usepackage{longtable}
\usepackage{array}
\usepackage{booktabs}
\usepackage{hyperref}
\usepackage{xcolor}
\geometry{left=2.0cm,right=2.0cm,top=2.0cm,bottom=2.0cm}

\title{Cross-View Geo-Localization\\Tam Proje Yapısı Özeti (Ayrıntılı)}
\author{Otomatik Envanter ve Mimari Özeti}
\date{4 Nisan 2026}

\begin{document}
\maketitle

\section*{Kapsam ve Amaç}
Bu doküman, çalışma alanındaki \textbf{tüm görünen dosya yapısını} teknik olarak özetler:
\begin{itemize}
    \item Mimari katmanlar (model, eğitim, test, görselleştirme, yardımcı modüller),
    \item Her ana dosya grubunun görevi,
    \item Eğitim/değerlendirme akışlarının haritalanması,
    \item Ek bölümde \textbf{.git hariç tam dosya listesi}.
\end{itemize}

\section{Depo Profili (Anlık Envanter)}
\begin{itemize}
    \item Toplam dosya (\texttt{.git} hariç): \textbf{102}
    \item Toplam klasör (\texttt{.git} hariç): \textbf{25}
    \item Python dosyası: \textbf{40}
    \item Markdown dosyası: \textbf{9}
    \item LaTeX dosyası: \textbf{2}
    \item Notebook dosyası: \textbf{2}
    \item Görsel varlık: \textbf{31} (png/jpg)
\end{itemize}

\subsection*{Dosya Türü Dağılımı}
\begin{center}
\begin{tabular}{lr}
\toprule
Uzantı/Tür & Adet \\
\midrule
py & 40 \\
md & 9 \\
tex & 2 \\
ipynb & 2 \\
sh & 3 \\
cpp & 2 \\
cu & 2 \\
txt & 5 \\
png & 27 \\
jpg & 4 \\
html & 1 \\
m & 1 \\
.gitignore/.gitkeep/.DS\_Store vb. & 4 \\
\bottomrule
\end{tabular}
\end{center}

\section{Sistem Mimarisi (Fonksiyonel Katmanlar)}
\subsection{1) Model Katmanı}
\begin{itemize}
    \item \texttt{model.py}: Temel ağ bileşenleri (\texttt{GeM}, \texttt{ClassBlock}, \texttt{ft\_net}, \texttt{two\_view\_net}, \texttt{three\_view\_net}).
    \item \texttt{model\_rgbd.py}: RGBD varyantı (\texttt{convert\_conv1\_to\_4channel}, \texttt{ft\_net\_rgbd}, \texttt{two\_view\_net\_rgbd}).
    \item \texttt{circle\_loss.py}: Metric learning odaklı \texttt{CircleLoss} implementasyonu.
\end{itemize}

\subsection{2) Veri Katmanı}
\begin{itemize}
    \item \texttt{folder.py}: torchvision benzeri veri klasörü okuma altyapısı.
    \item \texttt{dataset\_rgbd.py}: CVUSA/CVUSA-RGBD veri seti sınıfları ve RGB+Depth birleştirme.
    \item \texttt{prepare\_cvusa.py}, \texttt{prepare\_limit\_view.py}: veri hazırlama/ön işleme scriptleri.
\end{itemize}

\subsection{3) Eğitim Katmanı}
\begin{itemize}
    \item Ana eğitim: \texttt{train.py}, \texttt{train\_cvusa.py}
    \item Deneysel varyantlar: \texttt{train\_no\_street.py}, \texttt{train\_1\_sample.py}, \texttt{train\_3\_sample.py}, \texttt{train\_9\_sample.py}, \texttt{train\_18\_sample.py}, \texttt{train\_27\_sample.py}
    \item Çalıştırma scriptleri: \texttt{train\_baseline\_rgb.sh}, \texttt{train\_rgbd.sh}
\end{itemize}

\subsection{4) Test ve Değerlendirme Katmanı}
\begin{itemize}
    \item Genel test: \texttt{test.py}, \texttt{test\_train.py}
    \item Dataset/ölçek özel testler: \texttt{test\_cvusa.py}, \texttt{test\_4K.py}, \texttt{test\_160k.py}, \texttt{test\_ImageNet.py}, \texttt{test\_Place365.py}
    \item Metrik değerlendirme: \texttt{evaluate\_gpu.py}
\end{itemize}

\subsection{5) Görselleştirme ve Analiz Katmanı}
\begin{itemize}
    \item Demo: \texttt{demo.py}, \texttt{demo\_4K.py}
    \item Grad-CAM: \texttt{gradcam\_visualization.py}, \texttt{compare\_gradcam.py}
    \item Model karşılaştırma: \texttt{compare\_models.py}
    \item RGBD iyileşme analizi: \texttt{find\_rgb\_improvements.py}
    \item Veri gösterimi: \texttt{show\_data.py}
\end{itemize}

\subsection{6) Re-ranking (GPU GNN) Alt Sistemi}
\begin{itemize}
    \item \texttt{GPU-Re-Ranking/gnn\_reranking.py}: GNN tabanlı yeniden sıralama.
    \item \texttt{GPU-Re-Ranking/evaluate\_rerank\_gpu.py}: performans ve metrik değerlendirme.
    \item \texttt{extension/}: CUDA/C++ çekirdekleri (adjacency ve propagation).
\end{itemize}

\section{Ana Dosya Grupları ve Sorumluluk Matrisi}
\begin{longtable}{|p{0.30\textwidth}|p{0.64\textwidth}|}
\hline
\textbf{Dosya(lar)} & \textbf{Görev / Teknik Not} \\
\hline
\endfirsthead
\hline
\textbf{Dosya(lar)} & \textbf{Görev / Teknik Not} \\
\hline
\endhead

\texttt{README.md} & Projenin ana açıklaması; University-1652 görev tanımı, eğitim/test örnek komutları, haberler ve referanslar. \\
\hline
\texttt{README\_RGBD.md} & MiDaS depth + RGBD mimarisinin entegrasyonu, 4-kanal giriş mantığı, beklenen kazanımlar. \\
\hline
\texttt{README\_GRADCAM.md} & Grad-CAM/Grad-CAM++ kullanım kılavuzu, karşılaştırmalı analiz ve metrik yorumlama. \\
\hline
\texttt{TRAINING\_GUIDE.md} & RGB vs RGBD eğitim akışları, önerilen hiperparametreler, beklenen metrik farkları. \\
\hline
\texttt{technical\_documentation.tex} & Teorik arka plan (GeM, Circle Loss, RGBD adaptasyonu), pipeline ve matematiksel formülasyonlar. \\
\hline
\texttt{model.py} & Backbone + pooling + classifier modülleri; iki/üç görünüm ağ tanımları. \\
\hline
\texttt{model\_rgbd.py} & İlk konvolüsyon katmanını 4 kanala dönüştüren ve RGBD dalını tanımlayan model dosyası. \\
\hline
\texttt{dataset\_rgbd.py} & \texttt{CVUSADataset}, \texttt{CVUSARGBDDataset}, \texttt{RGBDSatelliteDataset}, \texttt{MixedRGBDDataset}. \\
\hline
\texttt{train.py} & Çok-görüşlü ana eğitim scripti; scheduler, optimizer, mixed precision ve eğri çizimi. \\
\hline
\texttt{train\_cvusa.py} & CVUSA odaklı eğitim scripti (RGB ve opsiyonel RGBD yolları). \\
\hline
\texttt{test.py} & Genel feature çıkarımı ve retrieval değerlendirmesi. \\
\hline
\texttt{test\_cvusa.py} & CVUSA test akışı; \texttt{RGBDTransform} ile derinlik kanalının entegrasyonu. \\
\hline
\texttt{gradcam\_visualization.py} & \texttt{GradCAM}, \texttt{GradCAMPlusPlus} sınıfları ve tekli/batch CAM üretimi. \\
\hline
\texttt{compare\_gradcam.py} & RGB vs RGBD dikkat haritası fark analizi (entropy, peak, gini vb.). \\
\hline
\texttt{compare\_models.py} & Parametre sayımı, katman bazlı karşılaştırma ve raporlama. \\
\hline
\texttt{evaluate\_gpu.py} & mAP ve retrieval metriği hesaplama. \\
\hline
\texttt{autoaugment.py}, \texttt{random\_erasing.py} & Veri artırma stratejileri. \\
\hline
\texttt{utils.py} & Model kaydet/yükle, sınıf dengeleme ağırlıkları, EMA benzeri yardımcılar. \\
\hline
\texttt{folder.py} & Veri seti dizin okuma altyapısı (ImageFolder türevi). \\
\hline
\texttt{GPU-Re-Ranking/*} & GNN tabanlı post-processing ile gerçek-zamanlıya yakın yeniden sıralama. \\
\hline
\texttt{docs/*}, \texttt{tutorial/*}, \texttt{State-of-the-art/*} & Proje web dokümanı, öğretici materyal, literatür/benchmark notları. \\
\hline
\end{longtable}

\section{Eğitim/Test İş Akışı Haritası}
\subsection*{Standart University-1652 Akışı}
\begin{enumerate}
    \item Veri hazırlığı ve klasör doğrulama.
    \item Model seçimi: RGB (\texttt{model.py}) veya RGBD (\texttt{model\_rgbd.py}).
    \item Eğitim: \texttt{train.py} (veya dataset özel \texttt{train\_cvusa.py}).
    \item Test: \texttt{test.py} / \texttt{test\_cvusa.py}.
    \item Sonuç görselleştirme: \texttt{demo.py}, Grad-CAM dosyaları.
    \item Opsiyonel post-processing: \texttt{GPU-Re-Ranking}.
\end{enumerate}

\subsection*{RGBD Akışı (CVUSA odaklı)}
\begin{enumerate}
    \item MiDaS ile depth üretimi (harici adım/Colab pipeline).
    \item \texttt{dataset\_rgbd.py} ile RGB+Depth yükleme.
    \item \texttt{train\_cvusa.py --use\_rgbd} ile eğitim.
    \item \texttt{test\_cvusa.py --use\_rgbd} ile ölçüm.
    \item \texttt{compare\_models.py} + \texttt{compare\_gradcam.py} ile analiz.
\end{enumerate}

\section{Alt Klasör Bazlı Ayrıntılar}
\subsection*{docs/}
\begin{itemize}
    \item Statik tanıtım sayfası: \texttt{index.html}
    \item Görsel ve referans varlıkları: \texttt{index\_files/}
\end{itemize}

\subsection*{GPU-Re-Ranking/}
\begin{itemize}
    \item Python API + CLI değerlendirme scriptleri
    \item CUDA extension kaynakları (\texttt{.cpp}/\texttt{.cu})
    \item \texttt{make.sh} ile derleme
\end{itemize}

\subsection*{SIFT/}
\begin{itemize}
    \item MATLAB demo scripti ve örnek görseller
\end{itemize}

\subsection*{image\_4K/}
\begin{itemize}
    \item 4K demo/test örnek query görüntüleri (00--11 klasörleri ve ek PNG'ler)
\end{itemize}

\subsection*{model/}
\begin{itemize}
    \item Eğitilmiş ağırlıklar için saklama noktası (şu an \texttt{.gitkeep})
\end{itemize}

\section{Ek A: .git Hariç Tam Dosya Listesi (102 Dosya)}
\small
\begin{verbatim}
.gitignore
COLAB_SATELLITE_DRONE.ipynb
GPU-Re-Ranking/README.md
GPU-Re-Ranking/evaluate_rerank_gpu.py
GPU-Re-Ranking/extension/adjacency_matrix/build_adjacency_matrix.cpp
GPU-Re-Ranking/extension/adjacency_matrix/build_adjacency_matrix_kernel.cu
GPU-Re-Ranking/extension/adjacency_matrix/setup.py
GPU-Re-Ranking/extension/make.sh
GPU-Re-Ranking/extension/propagation/gnn_propagate.cpp
GPU-Re-Ranking/extension/propagation/gnn_propagate_kernel.cu
GPU-Re-Ranking/extension/propagation/setup.py
GPU-Re-Ranking/gnn_reranking.py
GPU-Re-Ranking/utils.py
LICENSE
PROJECT_FILE_SUMMARY.tex
README.md
README_GRADCAM.md
README_RGBD.md
Request.md
SIFT/1.png
SIFT/2.png
SIFT/3.png
SIFT/README.md
SIFT/SIFT_Demo.png
SIFT/demo_sift.m
State-of-the-art/README.md
TRAINING_GUIDE.md
University1652_Colab.ipynb
autoaugment.py
circle_loss.py
colab_full_pipeline.py
compare_gradcam.py
compare_models.py
dataset_rgbd.py
demo.py
demo_4K.py
docs/index.html
docs/index_files/.DS_Store
docs/index_files/Motivation.png
docs/index_files/bibtex_cvpr2019.txt
docs/index_files/sample_drone.jpg
docs/index_files/sample_google.jpg
docs/index_files/sample_satellite.jpg
docs/index_files/sample_street.jpg
docs/index_files/top3.png
docs/index_files/youtube1.png
evaluate_gpu.py
find_rgb_improvements.py
folder.py
gallery_name.txt
gradcam_visualization.py
image_4K/00/query.png
image_4K/01/query.png
image_4K/02/query.png
image_4K/03/query.png
image_4K/04/query.png
image_4K/05/query.png
image_4K/06/query.png
image_4K/07/query.png
image_4K/08/query.png
image_4K/09/query.png
image_4K/1.png
image_4K/10/query.png
image_4K/11/query.png
image_4K/3.png
image_4K/6.png
image_4K/7.png
image_4K/8.png
image_4K/9.png
model.py
model/.gitkeep
model_rgbd.py
new_name_list.txt
prepare_cvusa.py
prepare_limit_view.py
query_name.txt
random_erasing.py
requirement.txt
show.png
show_data.py
technical_documentation.tex
test.py
test_160k.py
test_4K.py
test_ImageNet.py
test_Place365.py
test_cvusa.py
test_train.py
tool/clear_model.py
train.py
train_18_sample.py
train_1_sample.py
train_27_sample.py
train_3_sample.py
train_9_sample.py
train_baseline_rgb.sh
train_cvusa.py
train_no_street.py
train_rgbd.sh
tutorial/README.md
tutorial/RKNet.png
utils.py
\end{verbatim}
\normalsize

\section*{Son Not}
Bu özet, proje kökündeki yürütülebilir scriptleri, dokümantasyonları, görselleştirme araçlarını, CUDA uzantılarını ve veri/çıktı varlıklarını tek belgede birleştirir. Proje içinde fonksiyonel ayrışma nettir: \emph{modelleme} $\rightarrow$ \emph{eğitim} $\rightarrow$ \emph{test} $\rightarrow$ \emph{analiz} $\rightarrow$ \emph{re-ranking}.

\end{document}
